\documentclass[12pt,reqno]{amsart}
\usepackage[margin=1.1in]{geometry}

\usepackage{amsthm}

\usepackage{/Users/taylor/tex/pkg/general}
\usepackage{/Users/taylor/tex/pkg/fonts}
\usepackage{/Users/taylor/tex/pkg/colors}
\usepackage{/Users/taylor/tex/pkg/unicode}
\usepackage{/Users/taylor/tex/pkg/section}
\usepackage{/Users/taylor/tex/pkg/theorem}

\renewcommand{\.}{\mspace{1.5mu}}
\newcommand{\?}{\mspace{-1.5mu}}
\newcommand{\mdot}{\?\cdot\?}


\DUC{03BC}{\mu} 
\DUC{1D4D0}{\mathscr{A}} 

\renewcommand{\labelitemi}{\raisebox{0.5\height}{\scalebox{0.65}{\textbullet}}}
\renewcommand{\labelitemii}{\raisebox{0.5\height}{\scalebox{0.65}{$\circ$}}}

% \numberwithin{equation}{section}
% \numberwithin{figure}{section}
% 
% \setlist[enumerate]{
%     leftmargin=*,
%     align=left,
%     label=\textbf{\thesection.\arabic*},
%     labelindent=0pt
% }

\newcommand{\Impl}{\Longrightarrow}

\usepackage{eucal}
\usepackage[scaled=1.15,mathscr]{urwchancal} % Changes mathscr font


% Probability things
% ===========================
\newcommand{\E}{E}
\renewcommand{\P}{P} % default \P is paragraph symbol
\newcommand{\V}{V}
\DeclareMathOperator{\Cov}{\mathrm{Cov}}
\DeclareMathOperator{\Corr}{\mathrm{Corr}}
\DeclareMathOperator{\PV}{\mathrm{PV}}


\renewcommand{\c}{\mathbf{c}} % \c{c} gives "cedilla" diacritic

\newcommand{\vocab}[1]{\textbf{{\color{Blue}#1}}}
\newcommand{\lh}{\left}
\newcommand{\xS}{\Sigma}



%%%%%%%%%%%%%%%%%%%%%%%%%%%%%%%%%%%%%%%%%%%%%%%%%
% DOCUMENT BEGIN
%%%%%%%%%%%%%%%%%%%%%%%%%%%%%%%%%%%%%%%%%%%%%%%%%
\begin{document}

\tableofcontents





%%%%%%%%%%%%%%%%%%%%%%%%%%%%%%%%%%%%%%%%%%%%%%%%%
\section{Summary}
%%%%%%%%%%%%%%%%%%%%%%%%%%%%%%%%%%%%%%%%%%%%%%%%%


\begin{definition}
    We are discussing a Bond with the following terms:
    \begin{itemize}
        \item $A =$ \emph{par value}, the amount paid at maturity;
        \item $n =$ number of compounding periods \underline{per year};
        \item $T$ is the number of years \emph{remaining} to maturity;
            \begin{itemize}[itemindent=0pt]
                \item $N:=nT$ is the \emph{total} number of compounding periods. 
            \end{itemize}
        \item $Y$ is the \underline{annual} ``desired/prevailing\footnotemark'' rate of return (a.k.a.~the \emph{yield to maturity}) %
        \footnotetext{We can assume this is the \emph{risk-free rate}.} %
            \begin{itemize} 
                \item $y := Y/n$ is the per-period rate of return. 
            \end{itemize}
        \item $R$ is the \emph{annual} coupon rate; 
            \begin{itemize} 
                \item $ρ := R/n$ is the \emph{per-period} coupon rate.
            \end{itemize}
    \end{itemize}
    Finally, we let $P := P(ρ, y, N)$ be the \emph{current} price of the bond.
\end{definition}

We may derive (see below) that
    \(
        P(ρ, y, N) = A + A\big(\tfrac{ρ}{y}-1\big)\big(1-(1+y)^{-N}\big).  
    \)


\vspace{3\baselineskip}
% ===============================================
\subsection{Derivation}

If $y = Y/n$ is the ``prevailing'' per-period rate for the class of Bond, then the per-period \emph{discount factor} is
    \[
        𝜈 := (1 + Y/n) = (1 + y).
    \]
With this, the present value of the coupon payments (each of size $Aρ$) and the present value of the face value are
    \[
        \PV(\text{coupons}) + \PV(\text{Par}) = 
        Aρ(𝜈^{-1} + \cdots + 𝜈^{-N}) + A𝜈^{-N} = Aρ\mdot\fr{1-𝜈^{-N}}{𝜈(1-𝜈^{-1})} + A𝜈^{-N}.
    \]
Simplifying and using the fact that $𝜈-1 = y$, we arrive at the formula
    \[
        P = A + A\Big(\frac{ρ}{y}-1\Big)\big(1-(1+y)^{-N}\big)  
    \]
% ===============================================
\subsection{Extra/Side Notes}

\begin{enumerate}
    \item In the context of Bonds, the coupon rate is taken to be \underline{constant} for the life of the Bond.
\end{enumerate}



%%%%%%%%%%%%%%%%%%%%%%%%%%%%%%%%%%%%%%%%%%%%%%%%%
% DOCUMENT END
%%%%%%%%%%%%%%%%%%%%%%%%%%%%%%%%%%%%%%%%%%%%%%%%%
\end{document}